\section{Il protocollo Signal: l'accordo delle chiavi}\label{sec:AccordoChiavi}

\subsection{Comunicare con chi è offline}\label{sec:offline}

Per cifrare una conversazione, due interlocutori devono prima condividere una chiave
segreta, e in protocolli come TLS questa viene negoziata in tempo reale. Nella
messaggistica, però, il destinatario spesso non è connesso. Signal risolve il problema
introducendo un server in cui ciascun utente deposita in anticipo alcune chiavi
pubbliche. Ogni chiave pubblica ha una corrispondente chiave privata, che non lascia
mai il dispositivo. Alice può così stabilire una chiave condivisa con Bob anche se Bob
è offline~\cite{x3dh,pqxdh}.

\subsection{Le chiavi in gioco}\label{sec:chiavi}

Il protocollo usa quattro tipi di chiavi. La chiave
di identità (IK) è a lungo termine e consente di riconoscere l'utente. La
\emph{signed prekey} (SPK) è firmata digitalmente con la IK, così che se ne possa
verificare la provenienza, e viene sostituita periodicamente. Le \emph{one-time prekey}
(OPK) sono chiavi di scorta, ciascuna usata una sola volta. La chiave effimera (EK),
infine, è generata da Alice a ogni esecuzione e poi scartata.

\subsection{Il funzionamento}\label{sec:x3dh}

Il protocollo, denominato X3DH (\emph{Extended Triple Diffie-Hellman}), si basa
sullo scambio Diffie-Hellman (DH), che combinando la chiave privata di una parte con
quella pubblica dell'altra produce lo stesso segreto su entrambi i lati, senza che
questo transiti in rete~\cite{x3dh}.

Nella prima fase Bob pubblica sul server la propria IK, la SPK con la relativa firma
e un lotto di OPK. Nella seconda Alice scarica queste chiavi, verifica la firma della
SPK (se non è valida, il protocollo si interrompe) e calcola quattro scambi DH:
\begin{align*}
  DH_1 &= DH(IK_A, SPK_B) & DH_2 &= DH(EK_A, IK_B)\\
  DH_3 &= DH(EK_A, SPK_B) & DH_4 &= DH(EK_A, OPK_B)
\end{align*}
Una funzione di derivazione (KDF) li combina in un'unica chiave di sessione
$SK = \mathrm{KDF}(DH_1 \| DH_2 \| DH_3 \| DH_4)$.
$DH_1$ e $DH_2$ usano le chiavi di identità e servono ad autenticare le due parti.
$DH_3$ e $DH_4$ usano chiavi temporanee che vengono poi cancellate, per cui anche chi
in futuro rubasse le chiavi di identità non potrebbe decifrare le sessioni passate
(\emph{forward secrecy})~\cite{x3dh}. Nella terza fase Bob riceve il primo
messaggio, contenente $IK_A$, $EK_A$ e gli identificativi delle prekey usate, ricalcola
SK e cancella la OPK consumata~\cite{pqxdh}.

Il protocollo garantisce inoltre la \emph{negabilità}. I messaggi non vengono firmati,
ma autenticati con la chiave di sessione, che conoscono sia Alice sia Bob. Bob quindi
non può dimostrare a terzi che un messaggio l'ha scritto Alice, perché avrebbe potuto
scriverlo lui stesso.

\subsection{Il ruolo del server}\label{sec:server}

Il server non può calcolare SK, ma può comunque interferire. Se trattenesse le OPK, la
forward secrecy dipenderebbe soltanto dalla durata della SPK. Più grave è il caso in
cui sostituisca le chiavi pubbliche di un utente, realizzando un attacco
\emph{man-in-the-middle}. Per escluderlo, gli utenti devono confrontare le chiavi di
identità tramite un canale esterno, come impronte numeriche o QR code; senza questa
verifica non vi è garanzia sull'identità dell'interlocutore (\cref{sec:ScambioMessaggi}).

\subsection{La minaccia quantistica: PQXDH}\label{sec:pqxdh}

Un calcolatore quantistico abbastanza potente potrebbe rompere DH. Chi registra oggi
il traffico cifrato potrebbe quindi decifrarlo in futuro (\emph{harvest now, decrypt
later}). Per questo nel 2023 Signal ha introdotto PQXDH~\cite{signalpqblog}, che
affianca a DH un meccanismo post-quantistico chiamato Kyber. Il segreto prodotto da
Kyber viene usato insieme ai DH per derivare SK, così un attaccante dovrebbe rompere
entrambi.

La protezione è però parziale, perché l'autenticazione resta vulnerabile al calcolo
quantistico: un attaccante attivo potrebbe impersonare gli utenti. Un'analisi formale
indipendente ha inoltre individuato alcune imprecisioni nella specifica, poi corrette da Signal in una nuova
versione~\cite{cryspenpqxdh}.

\subsection{L'adozione in WhatsApp}\label{sec:whatsapp}

Dal 2016 WhatsApp adotta il protocollo Signal per messaggi e chiamate. Al momento della
registrazione, il client invia al server IK, SPK con firma e un lotto di OPK, secondo lo
schema appena descritto~\cite{whatsappwhitepaper}. Il codice di WhatsApp però non è
pubblico, quindi non si può verificare che l'app faccia davvero quello che descrive il
white paper.
